\section{Relocation and Practicality Results}
\label{sec:results_relocation}

The relocation figures complement the radio-performance results by quantifying the physical reconfiguration burden associated with movable infrastructure. The scheduling outputs report how often the movable \glspl{ap} change position and how much displacement the resulting schedules induce.

\begin{figure*}[t]
    \centering
    \begin{subfigure}[t]{0.47\textwidth}
        \centering
        \begingroup
        \def\postprocessfigurewidth{\linewidth}
        \pgfplotsset{table/search path={\ScheduleFigureRoot}}
        \input{\ScheduleOverviewFigurePath}
        \endgroup
        \caption{Strategy-level overview of total relocation distance and move fraction.}
        \label{fig:schedule_overview}
    \end{subfigure}\hfill
    \begin{subfigure}[t]{0.47\textwidth}
        \centering
        \begingroup
        \def\postprocessfigurewidth{\linewidth}
        \pgfplotsset{table/search path={\ScheduleFigureRoot}}
        \input{\ScheduleHistogramFigurePath}
        \endgroup
        \caption{Distribution of relocation distances across observed \gls{ap} moves.}
        \label{fig:schedule_histogram}
    \end{subfigure}
    \caption{Relocation-cost figures used to interpret how much physical infrastructure movement is required by the compared placement strategies.}
    \label{fig:schedule_pair}
\end{figure*}

Together, the panels in \cref{fig:schedule_pair} characterize the relocation activity associated with the evaluated placement strategies through aggregate movement summaries and relocation-distance distributions.
\FloatBarrier
